% Part: computability
% Chapter: computability-theory
% Section: computable-sets

\documentclass[../../../include/open-logic-section]{subfiles}

\begin{document}

\olfileid{cmp}{thy}{cps}
\olsection{Himpunan Komputabel}

Pangerten komputabilitas bisa kita jembarake
saka fungsi komputabel menyang himpunan komputabel:

\begin{defn}
  Tetepna $S$ minangka himpunan wilangan asli.
  Mula $S$ \emph{komputabel} yen lan mung yen
  fungsi karakteristike~$\Char{S}$ komputabel.
  Kanthi tembung liya, $S$~komputabel yen lan
  mung yen fungsi
\[
\Char{S}(x) =
\begin{cases}
1 & \text{yen $x \in S$} \\
0 & \text{yen ora}
\end{cases}
\]
komputabel. Semono uga, relasi
$R(x_0, \dots, x_{k-1})$ komputabel yen lan
mung yen fungsi karakteristike komputabel.

Himpunan lan relasi komputabel uga diarani
\emph{bisa diputusake}.
\end{defn}

\begin{explain}
Gatekna yen saiki ana sawetara pangerten
komputabilitas: kanggo fungsi parsial, fungsi,
lan himpunan. Aja nganti bingung!{} \iftag{TMs}{Mesin
  Turing kang ngitung fungsi parsial ngasilake
  output fungsi iku kanggo nilai input nalika
  fungsi kasebut ditegesi; mesin Turing kang
  ngitung sawijining himpunan ngasilake $1$
  utawa~$0$ sawisé mutusake apa nilai input
  kalebu himpunan iku.}{}
\end{explain}

\end{document}
